\documentclass[a4paper,12pt]{article}
\usepackage{amsmath, amssymb}
\usepackage[margin=1in]{geometry}

\begin{document}

\begin{center}
  {\Large \textbf{KUIS 3}}
\end{center}

\vspace{0.5cm}

\begin{tabular}{ll}
  Mata Kuliah    & : Kalkulus 2                               \\
  Hari, tanggal  & : Senin, 08 April 2025                     \\
  Sifat          & : Open Cheat Sheet                         \\
  Dosen Pengampu & : Alvian Alif Hidayatullah, S.Mat., M.Mat. \\
\end{tabular}

\vspace{0.5cm}
\parindent=0pt
\rule{\textwidth}{2pt}
\vspace{0.1cm}


Diberikan 5 soal dengan bobot nilai yang sama dan boleh dikerjakan tidak berurutan.

Tuliskan: Nama, NRP, dan Nomor Kelas pada lembar jawaban Anda.

\vspace{0.5cm}

\begin{center}
  \textnormal{DILARANG MELAKUKAN TINDAK KECURANGAN WAKTU UJIAN} \\
  ``Setiap tindak kecurangan akan mendapatkan sanksi akademik.''
\end{center}

\vspace{0.1cm}
\rule{\textwidth}{2pt}
\vspace{0.5cm}

\begin{enumerate}
  \item Hitung integral berikut dengan integral parsial:
        \[
          \int e^{7x}x^2\,dx.
        \]

  \item Dengan menggunakan rumus reduksi, hitung integral berikut:

        \hspace*{1.4em}(a)
        \[
          \int \cos^6(2x)\,\sin^5(2x)\,dx.
        \]

        \hspace*{1.4em}(b)
        \[
          \int \tan^6\theta\,\sec^3\theta\,d\theta.
        \]

  \item Dengan integrasi pecahan parsial, hitunglah integral
        \[
          \int \frac{3x^2-x}{(x-2)(x^2+1)}\,dx
        \]

  \item Dengan menggunakan substitusi Trigonometri, hitunglah Integral dari
        \[
          \int x^2\sqrt{x^2-16}\,dx
        \]

  \item Dengan substitusi bentuk pangkat rasional, hitunglah
        \[
          \int \frac{t^{3/2}}{t+t^{4/3}}\,dt
        \]

\end{enumerate}

\end{document}
